\documentclass[a4paper]{article}

%------------------------------------------------
% الحزم
%------------------------------------------------

\usepackage{enumitem}
\usepackage[margin=2.5cm]{geometry}
\usepackage{amsmath,amsfonts,amssymb}
\usepackage[charter]{mathdesign}
\usepackage{graphicx}
\usepackage{array}
\usepackage{fancyhdr}
\usepackage{polyglossia}

%------------------------------------------------
% اللغات والخطوط
%------------------------------------------------

\setdefaultlanguage[
    calendar=gregorian,
    numerals=maghrib
]{arabic}

\setotherlanguage{english}

% إذا لم تكن الخطوط مثبتة في الجهاز، تُؤخذ من المجلد ../fonts/
\IfFontExistsTF{Amiri}{%
    \newfontfamily\arabicfont[Script=Arabic, Scale=1.15]{Amiri}
}{%
    \newfontfamily\arabicfont[
        Script=Arabic, Scale=1.15,
        Path=../fonts/, Extension=.ttf,
        UprightFont=*-Regular, BoldFont=*-Bold,
        ItalicFont=*-Italic, BoldItalicFont=*-BoldItalic
    ]{Amiri}
}

\IfFontExistsTF{AlBattar}{%
    \newfontfamily\arabicfontsf[Script=Arabic, Scale=1.2]{AlBattar}
    % خط العناوين (يعمل حتى بعد استعمال \textenglish)
    \newfontfamily\khatt[Script=Arabic, Scale=1.2]{AlBattar}
}{%
    \newfontfamily\arabicfontsf[
        Script=Arabic, Scale=1.2, Path=../fonts/,
        BoldFont=ae_AlBattar.ttf
    ]{ae_AlBattar.ttf}
    % خط العناوين (يعمل حتى بعد استعمال \textenglish)
    \newfontfamily\khatt[
        Script=Arabic, Scale=1.2, Path=../fonts/,
        BoldFont=ae_AlBattar.ttf
    ]{ae_AlBattar.ttf}
}

%------------------------------------------------
% إعداد الصفحة
%------------------------------------------------

\renewcommand{\footrulewidth}{1pt}
\renewcommand{\headrulewidth}{1pt}

\rhead{\khatt المدرسة العليا للأساتذة ببوزريعة}
\lhead{\khatt ملحقة بومرداس}

\cfoot{}
\lfoot{\khatt الأستاذ: د. بن عبد الرزاق يوسف}

\parindent 0pt
\setlength{\headsep}{15pt}
\setlength{\parskip}{3pt}

\pagestyle{fancy}

\setlist[enumerate,1]{
    label=\arabic*),
    itemsep=4pt,
    topsep=3pt
}

\setlist[itemize]{
    itemsep=2pt,
    topsep=2pt
}

\arraycolsep=3pt

%================================================
% بداية الوثيقة
%================================================

\begin{document}

%------------------------------------------------
% العنوان
%------------------------------------------------


\rule{\textwidth}{1pt}

\begin{center}

{\khatt\textbf{
الفصل 5: البحث عن المعلومات، تحليلها وتلخيصها
}}

\end{center}

\rule{\textwidth}{1pt}

%================================================
% الوضعية المحفزة
%================================================


%================================================
% النشاط الأول
%================================================

\section*{\khatt\underline{النشاط الأول: تحديد الحاجة إلى المعلومات}}

\begin{enumerate}

\item يريد طالب البحث عن الموضوع التالي:

\begin{center}

\textbf{
``تأثير الذكاء الاصطناعي على التعليم الجامعي''
}

\end{center}

حوّل الموضوع العام إلى ثلاثة أسئلة بحث أكثر تحديدًا.



\item ما الفرق بين:

\begin{center}

\textbf{موضوع البحث}

\qquad و \qquad

\textbf{سؤال البحث}؟

\end{center}

علّل إجابتك.

%------------------------------------------------

\item اقترح كلمات مفتاحية مناسبة للبحث عن الموضوع
السابق.

\begin{center}

\textenglish{AI}

\qquad
\textenglish{Artificial Intelligence}

\qquad
\textenglish{Education}

\end{center}

أضف أربع كلمات مفتاحية أخرى:

\begin{itemize}

    \item ............................................................

    \item ............................................................

    \item ............................................................

    \item ............................................................

\end{itemize}

%------------------------------------------------

\item قارن بين عبارات البحث التالية:

\begin{center}

\textbf{أ.} الذكاء الاصطناعي

\vspace{0.15cm}

\textbf{ب.} الذكاء الاصطناعي في التعليم الجامعي

\vspace{0.15cm}

\textbf{ج.} فوائد ومخاطر الذكاء الاصطناعي التوليدي
في التعليم الجامعي

\end{center}

أي عبارة أكثر تحديدًا؟ ولماذا؟

%------------------------------------------------

\item اكتب صيغة بحث أكثر دقة باستعمال كلمات مفتاحية
مناسبة.


\end{enumerate}

%================================================
% النشاط الثاني
%================================================

\section*{\khatt\underline{النشاط الثاني (واجب منزلي): البحث واختيار المصادر}}

\textbf{يُنجز هذا النشاط في المنزل ويُسلَّم في الحصة القادمة.}

\begin{enumerate}

\item ابحث عن معلومات حول:

\begin{center}

\textbf{
الذكاء الاصطناعي التوليدي في التعليم الجامعي
}

\end{center}

واختر ثلاثة مصادر مختلفة.

%------------------------------------------------

\item سجّل المعلومات الأساسية عن كل مصدر:

\begin{center}

\begin{tabular}{|
>{\centering\arraybackslash}m{1.2cm}|
>{\raggedleft\arraybackslash}m{4cm}|
>{\raggedleft\arraybackslash}m{3cm}|
>{\raggedleft\arraybackslash}m{3cm}|
>{\raggedleft\arraybackslash}m{3cm}|}

\hline

\textbf{رقم}
&
\textbf{عنوان المصدر}
&
\textbf{الكاتب / الجهة}
&
\textbf{التاريخ}
&
\textbf{نوع المصدر}
\\

\hline

1
&
&
&
&
\\[0.5cm]

\hline

2
&
&
&
&
\\[0.5cm]

\hline

3
&
&
&
&
\\[0.5cm]

\hline

\end{tabular}

\end{center}

%------------------------------------------------

\item صنّف المصادر التي عثرت عليها:

\begin{itemize}

    \item مصدر أكاديمي.
    \item موقع مؤسسة رسمية.
    \item مقال صحفي.
    \item مدونة أو موقع شخصي.
    \item محتوى مولد بواسطة الذكاء الاصطناعي.

\end{itemize}

%------------------------------------------------

\item لماذا لا ينبغي الاعتماد على أول نتيجة تظهر
في محرك البحث؟

%------------------------------------------------

\item اذكر ثلاثة معايير يمكن استعمالها لتقييم جودة
المصدر.

\end{enumerate}


\end{document}